\chapter{Differentiaalrekening}\label{ch:b2:diffcalc}

De rekening met twee veranderlijken uit het volume van
bachelorjaar 1 rijpt tot de differentiaalrekening van
afbeeldingen tussen genormeerde ruimten: de \emph{\omterm{def:b2:diffcalc:differential}{differentiaal}}
als beste lineaire benadering, de kettingregel in volle
algemeenheid, de symmetriestelling van Schwarz \emph{bewezen},
formules van Taylor, en de volledige tweede-ordeanalyse van
extrema. De stelling van de inverse functie, de kroon op de
theorie, wordt geformuleerd met haar bewijsstrategie --- een
vast punt volgens Banach --- expliciet gemaakt.

Overal is $U$ een \omterm{def:b2:metric:topology}{open} deelverzameling van $\R^n$ (of van een
genormeerde ruimte; het eindigdimensionale geval draagt alle
ideeën) en $f \colon U \to \R^m$.

\section{De differentiaal}

\begin{definition}\label{def:b2:diffcalc:differential}
$f$ heet \emph{differentieerbaar}\index{differentieerbaar} in $a$
wanneer er een (\omterm{def:b2:metric:continuity}{continue}) lineaire afbeelding $\dd f_a \colon \R^n
\to \R^m$ bestaat met
\[
f(a + h) = f(a) + \dd f_a(h) + o(\norm h)
\qquad (h \to 0).
\]
De afbeelding $\dd f_a$, de \emph{differentiaal}\index{differentiaal}
van $f$ in $a$, is uniek; haar matrix in de canonieke basissen is
de \emph{jacobimatrix}\index{jacobimatrix} $J_f(a) =
\bigl(\frac{\partial f_i}{\partial x_j}(a)\bigr)$.
Differentieerbaarheid impliceert \omterm{def:b2:metric:continuity}{continuïteit} en het bestaan van
alle richtingsafgeleiden $\dd f_a(v) = \lim_{t\to0}\frac{f(a + tv)
- f(a)}{t}$; het omgekeerde faalt (\cref{exo:b2:diffcalc:2}). Voor
$m = 1$ is $\dd f_a(h) = \langle \nabla f(a), h\rangle$: de
gradiënt uit bachelorjaar 1, nu begrepen als de vector die de
differentiaal voorstelt.
\end{definition}

\begin{theorem}[$C^1$-criterium]\label{thm:b2:diffcalc:c1}
Bestaan alle partiële afgeleiden van $f$ op $U$ en zijn zij
\omterm{def:b2:metric:continuity}{continu} in $a$, dan is $f$ \omterm{def:b2:diffcalc:differential}{differentieerbaar} in $a$. ``$C^1$ op
$U$'' --- \omterm{def:b2:metric:continuity}{continue} partiële afgeleiden --- impliceert dus overal
\omterm{def:b2:diffcalc:differential}{differentieerbaarheid}, met \omterm{def:b2:metric:continuity}{continue} \omterm{def:b2:diffcalc:differential}{differentiaal}.
\end{theorem}

\begin{proof}
Component voor component ($m = 1$ volstaat). Het bewijs uit
bachelorjaar 1 voor twee veranderlijken --- verplaats één
coördinaat tegelijk, pas op elk traject de
middelwaardestelling in één veranderlijke toe, en gebruik de
\omterm{def:b2:metric:continuity}{continuïteit} van de partiële afgeleiden in $a$ --- veralgemeent
woordelijk tot $n$ trajecten:
\[
f(a + h) - f(a) = \sum_{j=1}^{n} \bigl(f(a + h^{(j)}) - f(a +
h^{(j-1)})\bigr)
= \sum_j h_j\,\frac{\partial f}{\partial x_j}(\xi_j),
\]
waarbij $h^{(j)}$ de eerste $j$ coördinaten van $h$ bevriest en
$\xi_j$ op het $j$-de traject ligt; de \omterm{def:b2:metric:continuity}{continuïteit} maakt van elke
$\frac{\partial f}{\partial x_j}(\xi_j)$ de uitdrukking
$\frac{\partial f}{\partial x_j}(a) + o(1)$, en de fout is
$o(\norm h)$.
\end{proof}

\begin{theorem}[Kettingregel]\label{thm:b2:diffcalc:chainrule}
Is $f$ \omterm{def:b2:diffcalc:differential}{differentieerbaar} in $a$ en $g$ in $f(a)$, dan is $g \circ
f$ \omterm{def:b2:diffcalc:differential}{differentieerbaar} in $a$ met
\[
\dd(g \circ f)_a = \dd g_{f(a)} \circ \dd f_a ,
\qquad
J_{g\circ f}(a) = J_g\bigl(f(a)\bigr)\,J_f(a) :
\]
jacobimatrices vermenigvuldigen.\index{kettingregel}
\end{theorem}

\begin{proof}
Schrijf $f(a + h) = f(a) + \dd f_a(h) + \norm h\,\varepsilon_1(h)$
en $g(b + k) = g(b) + \dd g_b(k) + \norm k\,\varepsilon_2(k)$ met
$b = f(a)$ en $k = k(h) = \dd f_a(h) + \norm h \varepsilon_1(h)$.
Substitutie geeft
\[
g(f(a+h)) = g(b) + \dd g_b\bigl(\dd f_a(h)\bigr)
+ \norm h\,\dd g_b(\varepsilon_1(h)) + \norm{k}\,\varepsilon_2(k),
\]
en beide fouttermen zijn $o(\norm h)$: de eerste omdat $\dd g_b$
\omterm{def:b2:metric:continuity}{continu} is en $\varepsilon_1 \to 0$; de tweede omdat $\norm k \leq
C\norm h$ (begrensde lineaire afbeelding plus een kleine term) en
$\varepsilon_2(k) \to 0$ als $h \to 0$.
\end{proof}

\begin{example}[Radiale functies, eens en voor altijd]\label{ex:b2:diffcalc:radial}
Zij $r(x) = \norm x_2$ op $\R^n\setminus\{0\}$ en $f = g \circ r$
met $g$ een $C^1$-functie van één veranderlijke. Eerst is $r$
\omterm{def:b2:diffcalc:differential}{differentieerbaar} buiten $0$: uit $r^2 = \sum x_i^2$ volgt
\[
\frac{\partial r}{\partial x_i} = \frac{x_i}{r},
\qquad\text{dat wil zeggen}\qquad
\nabla r(x) = \frac{x}{\norm x} ,
\]
de radiale eenheidsvector (differentieer $r^2$ en deel --- of pas
de kettingregel toe op $\sqrt{\cdot}$). Vervolgens geeft de
kettingregel voor elke radiale functie
\[
\nabla f(x) = g'\bigl(\norm x\bigr)\,\frac{x}{\norm x} .
\]
Uitgewerkt geval: $g(r) = \frac1r$ levert $\nabla\frac{1}{\norm x}
= -\frac{x}{\norm x^3}$ op, het omgekeerd-kwadratische veld van
de zwaartekracht en de elektrostatica --- radiaal van richting,
met grootte $\frac{1}{\norm x^2}$. Het inzicht om te onthouden:
gradiënten van radiale functies zijn radiaal omdat de
niveauverzamelingen bollen zijn en de gradiënt loodrecht op de
niveauverzamelingen staat; in $x = 0$ daarentegen is $r$
\emph{niet} \omterm{def:b2:diffcalc:differential}{differentieerbaar} (geen enkele kandidaat-lineaire
afbeelding past bij $\norm h$ uit alle richtingen) --- gladde
radiale profielen hebben $g'(0) = 0$ nodig om de oorsprong
sierlijk te passeren.
\end{example}

\begin{theorem}[Middelwaardeongelijkheid]\label{thm:b2:diffcalc:mvi}
Zij $f$ \omterm{def:b2:diffcalc:differential}{differentieerbaar} op $U$ en zij het segment $\intcc{a}{b}
= \{a + t(b-a)\}$ bevat in $U$. Dan is
\[
\norm{f(b) - f(a)} \leq \norm{b - a}\,
\sup_{x \in \intcc{a}{b}} \vertiii{\dd f_x} .
\]
In het bijzonder is een differentieerbare afbeelding met
verdwijnende \omterm{def:b2:diffcalc:differential}{differentiaal} op een \emph{\omterm{def:b2:metric:connected}{samenhangende}} \omterm{def:b2:metric:topology}{open}
verzameling constant.
\end{theorem}

\begin{proof}
De functie $\varphi(t) = f(a + t(b-a))$ is \omterm{def:b2:diffcalc:differential}{differentieerbaar} op
$\intcc{0}{1}$ met $\varphi'(t) = \dd f_{a + t(b-a)}(b - a)$
(kettingregel), met \omterm{def:b2:nvs:norm}{norm} $\leq M\norm{b-a}$ waarbij $M$ het
getoonde supremum is. Voor $\R$-waardige $f$ besluit de
middelwaardeongelijkheid in één veranderlijke; voor vectorwaarden
past men haar toe op $t \mapsto \langle u, \varphi(t)\rangle$ met
$u$ de eenheidsvector langs $f(b) - f(a)$. Constantheid: lokaal
constant (segmenten in ballen) plus \omterm{def:b2:metric:connected}{samenhang} (de verzameling
waar $f$ een gegeven waarde aanneemt, is \omterm{def:b2:metric:topology}{open} en gesloten:
\cref{ch:b2:metric}).
\end{proof}

\begin{example}[Een lipschitzconstante uit de middelwaardeongelijkheid]\label{ex:b2:diffcalc:lipschitz}
Is $f(x, y) = \sin x\,\sin y$ lipschitz op $\R^2$, en met welke
constante? Haar gradiënt is $\nabla f = (\cos x\sin y,\ \sin x\cos
y)$, met kwadratische \omterm{def:b2:nvs:norm}{norm}
\[
\cos^2x\sin^2y + \sin^2x\cos^2y
\leq \sin^2 y + \cos^2y\cdot 1 = 1
\]
(begrens $\cos^2x$ en $\sin^2x$ afzonderlijk door $1$), dus
$\vertiii{\dd f_{(x,y)}} = \norm{\nabla f} \leq 1$ overal, en
\cref{thm:b2:diffcalc:mvi} op het segment tussen twee willekeurige
punten geeft
\[
\abs{f(b) - f(a)} \leq \norm{b - a}_2 :
\]
$f$ is $1$-lipschitz, en de constante is scherp (nabij de
oorsprong heeft $f(x, \tfrac\pi2) = \sin x$ helling $1$). Het
inzicht om te onthouden: de middelwaardeongelijkheid zet een
\omterm{def:b2:funcseq:def}{puntsgewijze} grens op de \omterm{def:b2:diffcalc:differential}{differentiaal} om in een globale
continuïteitsmodulus --- de standaardweg naar lipschitzschattingen
in elke dimensie, en de motor in \cref{exo:b2:diffcalc:12}.
\end{example}

\section{Tweede afgeleiden}

\begin{theorem}[Schwarz]\label{thm:b2:diffcalc:schwarz}
Is $f$ van klasse $C^2$ op $U$ (alle tweede partiële afgeleiden
bestaan en zijn \omterm{def:b2:metric:continuity}{continu}), dan geldt voor alle $i,
j$:\index{stelling van Schwarz}
\[
\frac{\partial^2 f}{\partial x_i\,\partial x_j}
= \frac{\partial^2 f}{\partial x_j\,\partial x_i} .
\]
\end{theorem}

\begin{proof}
Twee veranderlijken volstaan ($x = x_i$, $y = x_j$, de rest
bevroren). Beschouw het tweede verschil
\[
\Delta(h) = f(a + h, b + h) - f(a + h, b) - f(a, b + h) + f(a,b) .
\]
Houd $h$ vast en zet $\varphi(x) = f(x, b+h) - f(x, b)$: dan is
$\Delta(h) = \varphi(a + h) - \varphi(a)$, en twee toepassingen van
de middelwaardestelling geven
\[
\Delta(h) = h\,\varphi'(\xi)
= h\Bigl(\frac{\partial f}{\partial x}(\xi, b+h) -
\frac{\partial f}{\partial x}(\xi, b)\Bigr)
= h^2\,\frac{\partial^2 f}{\partial y\,\partial x}(\xi, \eta),
\]
met $\xi \in \intoo{a}{a+h}$ en $\eta \in \intoo{b}{b+h}$. Wegens
de \omterm{def:b2:metric:continuity}{continuïteit} is $\frac{\Delta(h)}{h^2} \to \frac{\partial^2
f}{\partial y\partial x}(a, b)$ als $h \to 0$. Dezelfde berekening
met verwisselde rollen van de veranderlijken (bevries eerst de
tweede veranderlijke) geeft $\frac{\Delta(h)}{h^2} \to
\frac{\partial^2 f}{\partial x \partial y}(a,b)$: de twee limieten
van dezelfde grootheid vallen samen.
\end{proof}

\begin{example}[Waarom $C^2$ nodig is: het tegenvoorbeeld van Peano]\label{ex:b2:diffcalc:peano}
Zij $f(x, y) = \dfrac{xy(x^2 - y^2)}{x^2 + y^2}$, $f(0,0) = 0$.
Buiten de oorsprong is $f$ van klasse $C^\infty$; in de oorsprong
bestaan alle eerste en tweede partiële afgeleiden, maar de gemengde
zijn verschillend. Reken langs de assen: $f(x, 0) = f(0, y) = 0$,
en voor $y \neq 0$ is
\[
\frac{\partial f}{\partial x}(0, y)
= \lim_{x\to0}\frac{f(x,y)}{x}
= \frac{y(0 - y^2)}{y^2} = -y ,
\qquad\text{en symmetrisch}\qquad
\frac{\partial f}{\partial y}(x, 0) = x .
\]
Bijgevolg is
\[
\frac{\partial^2 f}{\partial y\,\partial x}(0,0)
= \frac{\dd}{\dd y}\Bigl[\frac{\partial f}{\partial
x}(0,y)\Bigr]_{y=0} = -1,
\qquad
\frac{\partial^2 f}{\partial x\,\partial y}(0,0) = +1 :
\]
de twee gemengde partiële afgeleiden bestaan en verschillen. Geen
tegenspraak met \cref{thm:b2:diffcalc:schwarz}: de tweede partiële
afgeleiden van $f$ zijn niet \omterm{def:b2:metric:continuity}{continu} in $0$ (test langs $y = tx$).
Het inzicht om te onthouden: de stelling van Schwarz is een echte
stelling over \emph{\omterm{def:b2:metric:continuity}{continuïteit}}, geen formele identiteit --- en
de hypothese ``$C^2$'' in Taylor--Young hieronder doet echt werk.
\end{example}

\begin{theorem}[Taylor--Young van orde 2]\label{thm:b2:diffcalc:taylor}
Zij $f \colon U \to \R$ van klasse $C^2$ en $a \in U$. Dan geldt,
als $h \to 0$,
\[
f(a + h) = f(a) + \langle\nabla f(a), h\rangle
+ \frac12\, \langle H_a h,\, h\rangle + o\bigl(\norm h^2\bigr),
\]
waarbij $H_a = \bigl(\frac{\partial^2 f}{\partial x_i\partial
x_j}(a)\bigr)$ de (volgens Schwarz \omterm{def:b2:quadratic:adjoint}{symmetrische})
\emph{hessiaan}\index{hessiaan} is.
\end{theorem}

\begin{proof}
Pas de stelling van Taylor--Young in één veranderlijke (volume van
bachelorjaar 1) toe op $\varphi(t) = f(a + th)$ op $\intcc{0}{1}$:
volgens de kettingregel is $\varphi'(t) = \langle \nabla f(a + th),
h\rangle$ en $\varphi''(t) = \langle H_{a+th}h, h\rangle$, beide
\omterm{def:b2:metric:continuity}{continu} in $t$. Dan is $\varphi(1) = \varphi(0) + \varphi'(0) +
\frac12\varphi''(\theta)$ (Taylor--Lagrange) met $\theta \in
\intoo{0}{1}$, en de \omterm{def:b2:metric:continuity}{continuïteit} van de tweede partiële
afgeleiden zet $\varphi''(\theta)$ om in $\varphi''(0) +
o(1)\cdot\norm h^2$, uniform: de getoonde ontwikkeling.
\end{proof}

\begin{example}[Eén ontwikkeling, twee wegen]\label{ex:b2:diffcalc:taylortwoways}
Ontwikkel $f(x, y) = \eu^x\cos y$ in de oorsprong tot orde $2$.
\emph{Door samenstelling van ontwikkelingen in één veranderlijke:}
\[
\eu^x\cos y = \Bigl(1 + x + \frac{x^2}{2} +
o(x^2)\Bigr)\Bigl(1 - \frac{y^2}{2} + o(y^2)\Bigr)
= 1 + x + \frac{x^2 - y^2}{2} + o\bigl(\norm{(x,y)}^2\bigr) .
\]
\emph{Via partiële afgeleiden:} $f_x = \eu^x\cos y$, $f_y =
-\eu^x\sin y$, dus $\nabla f(0) = (1, 0)$; en $f_{xx} = f$,
$f_{yy} = -f$, $f_{xy} = -\eu^x\sin y$ geven $H_0 =
\operatorname{diag}(1, -1)$: \cref{thm:b2:diffcalc:taylor}
reproduceert $1 + x + \frac12(x^2 - y^2)$. De twee berekeningen
stemmen overeen, en de weg via samenstelling was sneller ---
helemaal geen tweede partiële afgeleiden. Het inzicht om te
onthouden: de oorsprong is \emph{geen} kritiek punt ($\nabla f \neq
0$), zodat er ondanks de onbepaalde hessiaan geen zadelpunt te
melden valt: de lineaire term regeert, en de tweede-ordetest
spreekt uitsluitend in kritieke punten.
\end{example}

\begin{theorem}[Tweede-ordetest voor extrema, bewezen]\label{thm:b2:diffcalc:extrema}
Zij $f$ van klasse $C^2$ nabij een kritiek punt $a$ ($\nabla f(a) =
0$), met hessiaan $H = H_a$.
\begin{enumerate}
  \item Is $H$ positief definiet, dan is $a$ een strikt lokaal
        minimum (negatief definiet: maximum).
  \item Heeft $H$ \omterm{def:b2:reduction:eigen}{eigenwaarden} van beide tekens, dan is $a$ een
        zadelpunt: geen extremum.
  \item Is $H$ singulier (en semidefiniet), dan geen besluit.
\end{enumerate}
De test ``$rt - s^2$'' uit bachelorjaar 1 is het geval $n = 2$:
$\det H = rt - s^2$, en het teken van het spoor leest men af van
$r$.
\end{theorem}

\begin{proof}
Volgens de spectraalstelling (\cref{thm:b2:quadratic:spectral})
wordt de \omterm{def:b2:quadratic:def}{kwadratische vorm} van $H$ ingeklemd tussen haar extreme
\omterm{def:b2:reduction:eigen}{eigenwaarden}: $\lambda_{\min}\norm h^2 \leq \langle Hh, h\rangle
\leq \lambda_{\max}\norm h^2$.

(1) Is $\lambda_{\min} > 0$, dan geeft Taylor--Young
\[
f(a + h) - f(a) \geq \frac{\lambda_{\min}}{2}\norm h^2 -
o(\norm h^2) > 0
\]
voor kleine $h \neq 0$: strikt lokaal minimum.

(2) Langs een \omterm{def:b2:reduction:eigen}{eigenvector} $v_+$ met $\lambda_+ > 0$ is $f(a + tv_+)
- f(a) = \frac{\lambda_+}{2}t^2 + o(t^2) > 0$ voor kleine $t$;
langs $v_-$ met $\lambda_- < 0$ is het verschil negatief: beide
tekens komen in elke omgeving voor.

(3) $f(x,y) = x^2 + y^4$, $x^2 - y^4$ en $x^2 + y^3$ delen in $0$
dezelfde singuliere semidefiniete hessiaan, met drie verschillende
gedragingen.
\end{proof}

\begin{example}[Een volledige klassering, globaal inbegrepen]\label{ex:b2:diffcalc:quarticrun}
Klasseer alle extrema van $f(x, y) = x^4 + y^4 - 4xy$ op $\R^2$.
\emph{Kritieke punten:} $\nabla f = (4x^3 - 4y,\ 4y^3 - 4x) = 0$
geeft $y = x^3$ en $x = y^3 = x^9$, dus $x(x^8 - 1) = 0$: de reële
oplossingen zijn $(0,0)$, $(1,1)$ en $(-1,-1)$. \emph{Hessianen:}
$H = \begin{pmatrix} 12x^2 & -4\\ -4 & 12y^2\end{pmatrix}$. In
$(\pm1, \pm1)$: $\begin{pmatrix} 12 & -4\\ -4 &
12\end{pmatrix}$, \omterm{def:b2:reduction:eigen}{eigenwaarden} $8$ en $16$: positief definiet,
strikte lokale minima met $f = -2$. In $(0,0)$: $\begin{pmatrix} 0
& -4\\ -4 & 0\end{pmatrix}$, \omterm{def:b2:reduction:eigen}{eigenwaarden} $\pm4$: een zadelpunt.
\emph{Globaliteit:} uit $2\abs{xy} \leq x^2 + y^2$ volgt
\[
f(x, y) \geq x^4 + y^4 - 2(x^2 + y^2)
= (x^2 - 1)^2 + (y^2 - 1)^2 + x^2 + y^2 - 2
\xrightarrow[\norm{(x,y)}\to\infty]{} +\infty :
\]
$f$ is coërcief en neemt dus een globaal minimum aan
(\omterm{def:b2:metric:compact}{compactheid} van de subniveauverzamelingen), noodzakelijk in een
kritiek punt: de waarde $-2$, in zowel $(1,1)$ als $(-1,-1)$, is
het globale minimum; een maximum is er niet ($f$ is naar boven
onbegrensd). Het inzicht om te onthouden: de lokale test klasseert
de kandidaten, maar alleen een groeiargument maakt van ``lokaal''
``globaal'' --- het tweestappenpatroon van elk optimalisatiebewijs
in dit boek.
\end{example}

\begin{method}[De extrema van $f \colon \R^n \to \R$ klasseren]\label{met:b2:diffcalc:extrema}
\begin{enumerate}
  \item Los $\nabla f = 0$ op (alle kritieke punten; op een gebied
        met rand behandel je de rand apart, zoals in
        \cref{exo:b2:diffcalc:7}).
  \item Bereken in elk kritiek punt de hessiaan en de tekens van
        haar \omterm{def:b2:reduction:eigen}{eigenwaarden} --- in dimensie $2$ volstaan $\det H$ en
        $\operatorname{tr} H$: $\det < 0$ zadelpunt; $\det > 0$
        extremum, van het type dat het teken van het spoor
        aangeeft; $\det = 0$: de test zwijgt, bestudeer $f$ langs
        krommen.
  \item Voeg voor globale uitspraken een argument van \omterm{def:b2:metric:compact}{compactheid}
        of coërciviteit toe ($f \to +\infty$ in het oneindige, of
        een \omterm{def:b2:metric:compact}{compacte} verzameling van nevenvoorwaarden), en
        vergelijk daarna de kritieke waarden.
\end{enumerate}
\end{method}

\section{De stelling van de inverse functie}

\begin{theorem}[Stelling van de inverse functie]\label{thm:b2:diffcalc:inverse}
Zij $f \colon U \to \R^n$ van klasse $C^1$ en $a \in U$ met $\dd
f_a$ \emph{inverteerbaar}. Dan bestaan er \omterm{def:b2:metric:topology}{open} omgevingen $V \ni a$
en $W \ni f(a)$ zodanig dat $f \colon V \to W$ een bijectie is met
$C^1$-inverse, en\index{stelling van de inverse functie}
\[
\dd (f^{-1})_{f(x)} = (\dd f_x)^{-1} \qquad (x \in V).
\]
\end{theorem}
\admitted

\begin{remark}[Waarom zij waar is: de strategie met het vaste punt]
Het oplossen van $f(x) = y$ nabij $a$ herschrijft zich als de
vastepuntsvergelijking $x = x + \dd f_a^{-1}\bigl(y - f(x)\bigr)
=: \Phi_y(x)$; de afbeelding $\Phi_y$ heeft \omterm{def:b2:diffcalc:differential}{differentiaal}
$\mathrm{id} - \dd f_a^{-1}\dd f_x$, klein nabij $a$ wegens de
\omterm{def:b2:metric:continuity}{continuïteit} van $\dd f$, zodat $\Phi_y$ een contractie is op een
kleine gesloten bal en de vastepuntsstelling van Banach
(\cref{thm:b2:metric:banach}) de unieke lokale oplossing $x =
f^{-1}(y)$ levert. \omterm{def:b2:metric:continuity}{Continuïteit} en \omterm{def:b2:diffcalc:differential}{differentieerbaarheid} van de
inverse volgen dan uit de schattingen in de contractie. De volledige
boekhouding wordt in bachelorjaar 3 uitgevoerd; de strategie --- en
de uitspraak --- worden vanaf nu vrij gebruikt. De bijbehorende
\emph{stelling van de impliciete functie} ($F(x, y) = 0$ oplossen
naar $y(x)$ wanneer $\frac{\partial F}{\partial y}$ inverteerbaar
is) volgt door de stelling toe te passen op $(x, y) \mapsto (x,
F(x,y))$.
\end{remark}

\begin{example}[Poolcoördinaten]\label{ex:b2:diffcalc:polar}
$\Phi(r, \theta) = (r\cos\theta, r\sin\theta)$ heeft \omterm{def:b2:diffcalc:differential}{jacobimatrix}
\[
J_\Phi = \begin{pmatrix} \cos\theta & -r\sin\theta\\ \sin\theta &
r\cos\theta \end{pmatrix},
\qquad \det J_\Phi = r :
\]
inverteerbaar voor $r \neq 0$, dus $\Phi$ is buiten de oorsprong
een lokaal $C^1$-diffeomorfisme --- de vergunning om ``over te
gaan op poolcoördinaten'', hernieuwd voor de meervoudige
integralen van \cref{ch:b2:multint}.
\end{example}

\begin{example}[Overal lokaal, nergens globaal]\label{ex:b2:diffcalc:localnotglobal}
Zij $f(x, y) = \bigl(\eu^x\cos y,\ \eu^x\sin y\bigr)$ op $\R^2$.
Haar \omterm{def:b2:diffcalc:differential}{jacobimatrix},
\[
J_f = \begin{pmatrix}
\eu^x\cos y & -\eu^x\sin y\\
\eu^x\sin y & \eu^x\cos y
\end{pmatrix},
\qquad
\det J_f = \eu^{2x} > 0 ,
\]
verdwijnt nooit: volgens \cref{thm:b2:diffcalc:inverse} is $f$ in
\emph{elk} punt van het vlak een lokaal $C^1$-diffeomorfisme. Toch
is $f$ verre van injectief: $f(x, y + 2\pi) = f(x, y)$, dus elke
waarde wordt oneindig vaak aangenomen; en surjectief is zij evenmin,
want $\norm{f(x, y)} = \eu^x > 0$ mist de oorsprong. Het inzicht om
te onthouden: de stelling van de inverse functie is onherleidbaar
\emph{lokaal} --- de inverteerbaarheid van elke $\dd f_a$ levert een
lappendeken van lokale inversen die niet tot één geheel hoeven samen
te vallen. (Wie complexe getallen kent, herkent $z \mapsto \eu^z$;
de lappendeken is de familie van logaritmetakken.) Vergelijk
\cref{exo:b2:diffcalc:12}, waar een kwantitatieve globale hypothese
wél één globale inverse afdwingt.
\end{example}

\begin{remark}[Klassieke valkuilen]
\emph{(i) Richtingsafgeleiden zijn goedkoop, differentialen niet:}
alle richtingsafgeleiden kunnen bestaan --- en zelfs niet lineair
van de richting afhangen --- zonder \omterm{def:b2:diffcalc:differential}{differentieerbaarheid}
(\cref{exo:b2:diffcalc:2}); alleen het $C^1$-criterium
(\cref{thm:b2:diffcalc:c1}) verheft partiële afgeleiden tot een
\omterm{def:b2:diffcalc:differential}{differentiaal}. \emph{(ii) Kritiek betekent niet extremaal:}
zadelpunten (\cref{ex:b2:diffcalc:quarticrun}) en het zwijgende
singuliere geval (\cref{thm:b2:diffcalc:extrema}~(3)) verbergen
zich beide achter $\nabla f = 0$. \emph{(iii) Geen vectorwaardige
middelwaardegelijkheid:} alleen de \emph{ongelijkheid} van
\cref{thm:b2:diffcalc:mvi} overleeft (\cref{exo:b2:diffcalc:9});
schrijf nooit $f(b) - f(a) = \dd f_c(b-a)$ voor $f$ met waarden in
$\R^m$, $m \geq 2$. \emph{(iv) Lokale inverteerbaarheid is geen
injectiviteit:} \cref{ex:b2:diffcalc:localnotglobal}. \emph{(v) De
gradiënt hoort bij het inproduct:} $\nabla f$ is de vector die $\dd
f_a$ voorstelt in een gekozen inproduct; verander het product (zoals
in het gewogen voorbeeld uit het hoofdstuk over \omterm{def:b2:quadratic:def}{kwadratische
vormen}) en de gradiënt draait mee, terwijl de \omterm{def:b2:diffcalc:differential}{differentiaal} --- het
intrinsieke object --- niet beweegt.
\end{remark}

\begin{remark}[Waar dit wordt gebruikt]
Alles stroomafwaarts van dit hoofdstuk is toegepaste
differentiaalrekening: het hoofdstuk over
differentiaalvergelijkingen lineariseert stromingen en gebruikt de
determinantformule van Liouville (bewezen in de weekendopgave van
dit hoofdstuk); de hoofdstukken over krommen en oppervlakken
bestuderen niveauverzamelingen en parametriseringen via de stelling
van de impliciete functie; meervoudige integralen wisselen van
veranderlijke via jacobimatrices. De weekendopgave ontwikkelt de
rekening \emph{op de ruimte van de matrices zelf} --- de
\omterm{def:b2:diffcalc:differential}{differentiaal} van de \omterm{def:b2:linalg:det}{determinant}, van de inverse, de
\omterm{ex:b2:nvs:matrixexp}{matrixexponentiaal}, en de orthogonale groep als gladde
niveauverzameling --- de schaduw in bachelorjaar 2 van wat het
volume van bachelorjaar 3 als variëteiten en lie-groepen
formaliseert.
\end{remark}

\section{Oefeningen}

\begin{exercise}[$\star$]\label{exo:b2:diffcalc:1}
Bereken de jacobimatrices van $f(x,y) = (x^2 - y^2,\, 2xy)$ en van
$\Phi(r,\theta,z) = (r\cos\theta, r\sin\theta, z)$; waar zijn de
differentialen inverteerbaar?
\end{exercise}

\begin{exercise}[$\star$]\label{exo:b2:diffcalc:2}
Zij $f(x,y) = \frac{x^3}{x^2 + y^2}$ ($f(0,0) = 0$). Bewijs dat
alle richtingsafgeleiden van $f$ in $0$ bestaan, maar dat $f$ niet
\omterm{def:b2:diffcalc:differential}{differentieerbaar} is in $0$ \emph{(de afbeelding $v \mapsto$
richtingsafgeleide is niet lineair)}.
\end{exercise}

\begin{exercise}[$\star$]\label{exo:b2:diffcalc:3}
Bepaal en klasseer de kritieke punten van $f(x, y) = x^3 + y^3 -
3xy$ met \cref{thm:b2:diffcalc:extrema}, en die van $g(x,y) = x^4 +
y^4 - 2(x - y)^2$.
\end{exercise}

\begin{exercise}[$\star\star$]\label{exo:b2:diffcalc:4}
Zij $f \colon \R^n \to \R$ van klasse $C^1$ en \emph{homogeen van
graad $p$}: $f(tx) = t^pf(x)$ voor $t > 0$. Bewijs de identiteit van
Euler
\[
\langle \nabla f(x), x\rangle = p\,f(x) ,
\]
en haar omkering voor $C^1$-functies op $\R^n\setminus\{0\}$.
\end{exercise}

\begin{exercise}[$\star\star$]\label{exo:b2:diffcalc:5}
Zij $A$ \omterm{def:b2:quadratic:adjoint}{symmetrisch} en $f(x) = \frac12\langle Ax, x\rangle -
\langle b, x\rangle$. Bereken $\nabla f$ en $H_f$; wanneer is $f$
convex? Toon aan, aangenomen dat $A$ positief definiet is, dat $f$
een uniek globaal minimum heeft in de oplossing van $Ax = b$ --- de
bestaansreden van het gradiëntafdalen.
\end{exercise}

\begin{exercise}[$\star\star$]\label{exo:b2:diffcalc:6}
(Multiplicator van Lagrange, één nevenvoorwaarde, met de hand
bewezen) Zijn $f, g$ van klasse $C^1$ op $\R^2$, en neem aan dat
$f$ in $a$ een lokaal extremum aanneemt op de niveauverzameling
$\{g = 0\}$, met $\nabla g(a) \neq 0$. Bewijs dat $\nabla f(a) =
\lambda\nabla g(a)$ voor een zekere $\lambda$. \emph{(Parametriseer
de niveauverzameling nabij $a$ met de stelling van de impliciete
functie en differentieer $t \mapsto f(\gamma(t))$.)} Toepassing: de
extrema van $f(x,y) = xy$ op de cirkel $x^2 + y^2 = 1$.
\end{exercise}

\begin{exercise}[$\star\star$]\label{exo:b2:diffcalc:7}
Bepaal de extrema van $f(x, y) = x^2 + y^2 - xy + x - y$ op $\R^2$,
en vervolgens haar maximum en minimum op de gesloten driehoek met
hoekpunten $(0,0)$, $(1,0)$, $(0,1)$ \emph{(inwendige kritieke
punten, dan de drie zijden, dan de hoekpunten)}.
\end{exercise}

\begin{exercise}[$\star\star\star$]\label{exo:b2:diffcalc:8}
Zij $f \colon \R^2 \to \R^2$, $f(x, y) = (x + y^2,\; y + x^2)$.
Toon aan dat $f$ nabij $0$ een lokaal diffeomorfisme is, bereken
$\dd(f^{-1})_{(0,0)}$, en bepaal de grootste $r$ waarvoor $\dd f$
inverteerbaar is op de bal $\norm{(x,y)}_2 < r$ \emph{(bereken
$\det J_f$)}.
\end{exercise}

\begin{exercise}[$\star\star\star$]\label{exo:b2:diffcalc:9}
(Rolle faalt, de middelwaarde overleeft) Geef een $f \colon \R \to
\R^2$ van klasse $C^1$ met $f(0) = f(2\pi)$ maar $f'(t) \neq 0$
voor alle $t$ (geen vectorwaardige Rolle). Ga vervolgens op je
voorbeeld de middelwaarde\emph{ongelijkheid} van
\cref{thm:b2:diffcalc:mvi} na.
\end{exercise}

\begin{exercise}[$\star$]\label{exo:b2:diffcalc:10}
Bereken de \omterm{def:b2:diffcalc:differential}{differentiaal} en de gradiënt van $f(x) = \norm x_2^2$ en
van $g(x) = \langle Ax, x\rangle$ op $\R^n$ ($A$ een vierkante
matrix, niet verondersteld \omterm{def:b2:quadratic:adjoint}{symmetrisch} te zijn), en de hessiaan van
elk. Voor welke $A$ is $g$ convex?
\end{exercise}

\begin{exercise}[$\star\star$]\label{exo:b2:diffcalc:11}
Zij $f \colon \R^n \to \R$ van klasse $C^2$. Bewijs dat $f$ convex
is dan en slechts dan als haar hessiaan $H_x$ in elke $x$ positief
semidefiniet is \emph{(herleid tot één veranderlijke: $t \mapsto
f(a + t(b-a))$; gebruik Taylor--Lagrange in één richting, en
evalueer voor de omkering $\varphi''$)}.
\end{exercise}

\begin{exercise}[$\star\star\star$]\label{exo:b2:diffcalc:12}
(Een globale inversiestelling) Zij $g \colon \R^n \to \R^n$ van
klasse $C^1$ met $\vertiii{\dd g_x} \leq k < 1$ voor alle $x$, en
$f = \mathrm{id} + g$.
\begin{enumerate}
  \item Toon aan dat $\norm{f(x) - f(y)} \geq (1 - k)\norm{x -
        y}$: $f$ is injectief, met \omterm{def:b2:metric:continuity}{continue} inverse op haar beeld.
  \item Toon aan dat voor elke $y \in \R^n$ de afbeelding $x
        \mapsto y - g(x)$ een contractie van de volledige ruimte
        $\R^n$ is, en besluit met de vastepuntsstelling van Banach
        (\cref{thm:b2:metric:banach}) dat $f$ surjectief is.
  \item Besluit dat $f$ een bijectie van $\R^n$ is met een
        $(1-k)^{-1}$-lipschitzinverse --- een \emph{globale}
        tegenhanger van \cref{thm:b2:diffcalc:inverse}, die
        daarentegen zuiver lokaal is.
\end{enumerate}
\end{exercise}

\section{Probleem: de rekening met matrices --- Jacobi, exponentiaal
en de orthogonale groep}

\begin{problem}\label{pb:b2:diffcalc:1}
De schoonste speeltuin voor differentiaalrekening is de ruimte
$\mathcal M_n(\R) \simeq \R^{n^2}$ zelf: haar natuurlijkste
afbeeldingen --- product, inverse, \omterm{def:b2:linalg:det}{determinant}, exponentiaal ---
hebben differentialen van opvallende elegantie. Deze opgave
berekent ze allemaal: de \emph{reeks van Neumann}, de
\omterm{def:b2:diffcalc:differential}{differentiaal} van de inverse, de \emph{formule van
Jacobi}\index{formule van Jacobi} voor de \omterm{def:b2:linalg:det}{determinant} met de
\emph{formule van Liouville} als dividend, de
\emph{matrixexponentiaal}\index{matrixexponentiaal} met $\det\eu^A
= \eu^{\operatorname{tr}A}$, en ten slotte de orthogonale groep
$O_n$ als gladde niveauverzameling met de antisymmetrische
matrices als raakruimte --- differentiaalmeetkunde in de kiem.
Overal is $\vertiii\cdot$ de \omterm{thm:b2:nvs:continuouslinear}{operatornorm} ondergeschikt aan
$\norm\cdot_2$, en $\langle X, Y\rangle =
\operatorname{tr}(X^{\mathsf T}Y)$ het inproduct van Frobenius.

\medskip
\noindent\textbf{Deel I --- De reeks van Neumann.}
\begin{enumerate}
  \item Bewijs de submultiplicativiteit $\vertiii{AB} \leq
        \vertiii A\,\vertiii B$ en leid af dat
        veeltermafbeeldingen van $A$ (matrixproducten,
        \omterm{def:b2:linalg:det}{determinant}, spoor) \omterm{def:b2:metric:continuity}{continu} zijn op $\mathcal M_n(\R)$.
  \item Toon voor $\vertiii X < 1$ aan dat $\sum_{k\geq0}X^k$
        \omterm{def:b2:series:def}{absoluut} convergeert in $\mathcal M_n(\R)$
        (\cref{thm:b2:nvs:absoluteconvergence}), dat haar som $(I
        - X)^{-1}$ is, en dat
        \[
        \vertiii{(I - X)^{-1}} \leq \frac{1}{1 - \vertiii X},
        \qquad
        (I - X)^{-1} = I + X + O\bigl(\vertiii X^2\bigr) .
        \]
  \item Leid af dat $GL_n(\R)$ \emph{\omterm{def:b2:metric:topology}{open}} is: is $A$
        inverteerbaar en $\vertiii H < \frac{1}{\vertiii{A^{-1}}}$,
        dan is $A + H$ inverteerbaar. Leid ook af dat $GL_n(\R)$
        \emph{dicht} ligt in $\mathcal M_n(\R)$ \emph{(verstoor $A$
        met $\varepsilon I$: $\det(A + \varepsilon I)$ is een
        veelterm in $\varepsilon$ die niet nul is)}.
  \item Toon aan dat de inversieafbeelding $\Phi(A) = A^{-1}$
        \omterm{def:b2:metric:continuity}{continu} is op $GL_n(\R)$.
\end{enumerate}

\medskip
\noindent\textbf{Deel II --- Eerste differentialen.}
\begin{enumerate}[resume]
  \item Toon aan dat de kwadrateringsafbeelding $A \mapsto A^2$
        \omterm{def:b2:diffcalc:differential}{differentieerbaar} is met \omterm{def:b2:diffcalc:differential}{differentiaal} $H \mapsto AH +
        HA$, en algemener dat $A \mapsto A^k$ de \omterm{def:b2:diffcalc:differential}{differentiaal} $H
        \mapsto \sum_{i=0}^{k-1} A^iHA^{k-1-i}$ heeft. Waarom mag
        men in het algemeen \emph{niet} $kA^{k-1}H$ schrijven?
  \item Bewijs dat $\Phi(A) = A^{-1}$ \omterm{def:b2:diffcalc:differential}{differentieerbaar} is op
        $GL_n(\R)$ met
        \[
        \dd\Phi_A(H) = -A^{-1}HA^{-1}
        \]
        \emph{(schrijf $(A + H)^{-1} = (I + A^{-1}H)^{-1}A^{-1}$
        en werk uit met vraag 2)}. Toets de formule aan het
        scalaire geval $n = 1$.
  \item Leid voor een $C^1$-kromme $t \mapsto A(t) \in GL_n(\R)$
        af dat $\bigl(A(t)^{-1}\bigr)' = -A^{-1}A'A^{-1}$, en
        ontwikkel $t \mapsto (I + tB)^{-1}$ tot op eerste orde in
        $t = 0$.
  \item Bereken de \omterm{def:b2:diffcalc:differential}{differentiaal} van $f(A) =
        \operatorname{tr}(A^k)$ en identificeer haar gradiënt voor
        het inproduct van Frobenius:
        \[
        \dd f_A(H) = k\operatorname{tr}\bigl(A^{k-1}H\bigr),
        \qquad
        \nabla f(A) = k\,\bigl(A^{k-1}\bigr)^{\mathsf T} .
        \]
  \item Dezelfde vragen voor $f(A) = \operatorname{tr}(A^{\mathsf
        T}A) = \norm A_F^2$: \omterm{def:b2:diffcalc:differential}{differentiaal}, gradiënt en de
        (constante) hessiaan; besluit dat $\norm\cdot_F^2$ strikt
        convex is.
\end{enumerate}

\medskip
\noindent\textbf{Deel III --- De formule van Jacobi.}
\begin{enumerate}[resume]
  \item Bewijs dat
        \[
        \det(I + H) = 1 + \operatorname{tr}H +
        O\bigl(\vertiii H^2\bigr)
        \]
        \emph{(werk $\det(e_1 + h_1, \dots, e_n + h_n)$ uit met
        de multilineariteit in de kolommen: termen met minstens
        twee $h$-kolommen zijn $O(\vertiii H^2)$)}: $\dd(\det)_I =
        \operatorname{tr}$.
  \item Leid voor inverteerbare $A$ af dat
        \[
        \dd(\det)_A(H) = \det(A)\,
        \operatorname{tr}\bigl(A^{-1}H\bigr) .
        \]
  \item Toon aan dat voor \emph{elke} $A$ (inverteerbaar of niet)
        geldt $\frac{\partial\det}{\partial a_{ij}}(A) = C_{ij}$,
        de cofactor op plaats $(i,j)$ \emph{(ontwikkeling van
        Laplace langs rij $i$)}, zodat met de geadjungeerde
        $\operatorname{adj}A = \operatorname{com}(A)^{\mathsf T}$:
        \[
        \dd(\det)_A(H) =
        \operatorname{tr}\bigl(\operatorname{adj}(A)\,H\bigr),
        \qquad
        \nabla(\det)(A) = \operatorname{com}(A) ,
        \]
        waarmee vraag 11 wordt teruggevonden als $A$ inverteerbaar
        is ($\operatorname{adj}A = \det(A)A^{-1}$). Dit is de
        \emph{formule van Jacobi}: $\bigl(\det A(t)\bigr)' =
        \operatorname{tr}\bigl(\operatorname{adj}(A(t))\,A'(t)\bigr)$.
  \item (Formule van Liouville) Zij $A(t)$ een $C^1$-kromme van
        matrices die voldoet aan de lineaire
        differentiaalvergelijking $A'(t) = M(t)A(t)$. Bewijs dat
        \[
        \bigl(\det A(t)\bigr)' =
        \operatorname{tr}\bigl(M(t)\bigr)\,\det A(t),
        \qquad\text{en dus}\qquad
        \det A(t) = \det A(0)\,
        \exp\Bigl(\int_0^t\operatorname{tr}M\Bigr)
        \]
        \emph{(gebruik $\operatorname{adj}(A)\,A = \det(A)I$ en de
        cyclische invariantie van het spoor)} --- de identiteit
        voor de wronskiaan die het hoofdstuk over
        differentiaalvergelijkingen voortdurend zal gebruiken.
  \item Toon aan dat $SL_n(\R) = \{\det = 1\}$ een gladde
        niveauverzameling is: in elke $A \in SL_n(\R)$ is de
        \omterm{def:b2:diffcalc:differential}{differentiaal} $\dd(\det)_A$ een \emph{surjectieve}
        lineaire afbeelding op $\R$ \emph{(evalueer haar in $H =
        \frac1nA$)}.
\end{enumerate}

\medskip
\noindent\textbf{Deel IV --- De \omterm{ex:b2:nvs:matrixexp}{matrixexponentiaal}.}
\begin{enumerate}[resume]
  \item Toon aan dat $\eu^A = \sum_{k\geq0}\frac{A^k}{k!}$ voor
        elke $A$ \omterm{def:b2:series:def}{absoluut} convergeert, normaal op elke bal, met
        $\vertiii{\eu^A} \leq \eu^{\vertiii A}$; en dat $\eu^A$
        \omterm{def:b2:metric:continuity}{continu} van $A$ afhangt.
  \item Bewijs dat uit $AB = BA$ volgt dat $\eu^{A+B} =
        \eu^A\eu^B$ \emph{(\omterm{thm:b2:series:fubini}{Cauchy-product}, geoorloofd wegens de
        \omterm{def:b2:series:def}{absolute convergentie})}; leid af dat $\eu^A$ altijd
        inverteerbaar is met inverse $\eu^{-A}$: $\exp$ beeldt
        $\mathcal M_n(\R)$ af in $GL_n(\R)$.
  \item Toon aan dat $t \mapsto \eu^{tA}$ van klasse $C^1$ is
        (zelfs $C^\infty$) met
        \[
        \frac{\dd}{\dd t}\,\eu^{tA} = A\,\eu^{tA} =
        \eu^{tA}A
        \]
        \emph{(differentieer de reeks term voor term op
        segmenten: de afgeleide reeks convergeert normaal)}.
  \item Bewijs de identiteit
        \[
        \det\bigl(\eu^{A}\bigr) = \eu^{\operatorname{tr}A}
        \]
        \emph{(pas de formule van Liouville, vraag 13, toe op
        $A(t) = \eu^{tA}$)}. Verstandscontroles: $n = 1$;
        nilpotente $A$; en de matrices met spoor nul belanden in
        $SL_n(\R)$.
  \item Toon aan dat $\eu^H = I + H + O(\vertiii H^2)$, zodat
        $\exp$ \omterm{def:b2:diffcalc:differential}{differentieerbaar} is in $0$ met $\dd(\exp)_0 =
        \mathrm{id}$; besluit met de stelling van de inverse
        functie (\cref{thm:b2:diffcalc:inverse}) dat $\exp$ een
        $C^1$-diffeomorfisme is van een omgeving van $0$ op een
        omgeving van $I$: elke matrix dicht bij de eenheidsmatrix
        heeft een logaritme.
  \item Toon aan dat $\exp$ de \omterm{def:b2:quadratic:adjoint}{symmetrische} matrices bijectief
        afbeeldt op de \omterm{def:b2:quadratic:adjoint}{symmetrische} \emph{positief definiete}
        matrices \emph{(diagonaliseer; de inverse is de spectrale
        logaritme)}.
\end{enumerate}

\medskip
\noindent\textbf{Deel V --- De orthogonale groep als
niveauverzameling.}
\begin{enumerate}[resume]
  \item Zij $F(A) = A^{\mathsf T}A$, van $\mathcal M_n(\R)$ naar
        de \omterm{def:b2:quadratic:adjoint}{symmetrische} matrices $S_n$. Bereken $\dd F_A(H) =
        A^{\mathsf T}H + H^{\mathsf T}A$ en toon aan dat deze
        \omterm{def:b2:diffcalc:differential}{differentiaal} in elke $A \in O_n = F^{-1}(I)$
        \emph{surjectief} is op $S_n$ \emph{(probeer voor gegeven
        $S \in S_n$ de matrix $H = \frac12 AS$)}: $O_n$ is een
        gladde niveauverzameling, van dimensie $n^2 -
        \frac{n(n+1)}2 = \frac{n(n-1)}2$.
  \item Toon aan dat elke $C^1$-kromme $A(t) \in O_n$ met $A(0) =
        I$ een \emph{antisymmetrische} snelheid $A'(0)$ heeft, en
        omgekeerd dat voor antisymmetrische $K$ de kromme
        $\eu^{tK}$ in $O_n$ blijft: de raakruimte van $O_n$ in $I$
        bestaat precies uit de antisymmetrische matrices.
  \item Toon aan dat $\det\eu^{K} = 1$ voor antisymmetrische $K$
        (vraag 18): de exponentiële kromme leeft in de
        rotatiegroep $SO_n$. Bereken haar \omterm{def:b2:metric:complete}{volledig} voor $n = 2$:
        bewijs met $J = \begin{pmatrix}0 & -1\\ 1 &
        0\end{pmatrix}$ dat
        \[
        \eu^{\theta J} = \begin{pmatrix} \cos\theta &
        -\sin\theta\\ \sin\theta & \cos\theta\end{pmatrix} :
        \]
        de \omterm{ex:b2:nvs:matrixexp}{matrixexponentiaal} \emph{is} de rotatie over $\theta$,
        en de reeksdefinities van cosinus en sinus duiken opnieuw
        op, binnen in een matrix.
  \item (De truc met de dichtheid) Breid met de dichtheid van
        $GL_n(\R)$ (vraag 3) en de \omterm{def:b2:metric:continuity}{continuïteit} de identiteit
        \[
        \operatorname{adj}(AB) =
        \operatorname{adj}(B)\operatorname{adj}(A)
        \]
        uit van inverteerbare tot alle matrices \emph{(voor
        inverteerbare $A, B$ zijn beide leden gelijk aan
        $\det(AB)(AB)^{-1}$; beide leden zijn veeltermen in de
        elementen)}.
  \item Synthese. In telkens één zin: (i) welke eerdere
        hoofdstukken de motor van elk deel leverden (\omterm{def:b2:metric:complete}{volledigheid}
        en genormeerde algebra's; de spectraalstelling; de
        stelling van de inverse functie); (ii) op welke formule
        uit deze opgave het hoofdstuk over
        differentiaalvergelijkingen zal steunen, en waar; (iii)
        wat $\dd(\det)_I = \operatorname{tr}$ en $\det\eu^A =
        \eu^{\operatorname{tr}A}$ zeggen over spoor en \omterm{def:b2:linalg:det}{determinant}
        als ``infinitesimaal en globaal volume''; (iv) wat het
        volume van bachelorjaar 3 van de vragen 21--23 maakt
        (lie-groepen en hun lie-algebra's).
\end{enumerate}
\end{problem}
